\section{Related Work}
\label{sec:related-work}

\textbf{Collective behavior in multi-agent language models.}
A growing literature asks what collective behavior emerges when language-model agents interact, and whether communication improves or degrades the resulting system. LLM populations can sustain cooperation in repeated social dilemmas and shared-resource problems \citep{piatti2024cooperate,fontana2025nicer,meulemans2026game}, form conventions and population-level biases \citep{ashery2025emergent}, and exhibit higher-order coordination that is not reducible to independent agent behavior \citep{riedl2026emergent}. At the same time, interaction is not uniformly beneficial: collectives can converge before adequately pooling distributed evidence \citep{hiddenbench2026}, become distorted by conformity, perceived expertise, and dominant speakers \citep{ko2026social}, exhibit free-riding or collusive behavior \citep{guzman2025corrupted,nakamura2026colosseum}, and suffer broader failures of miscoordination and brittle social epistemics \citep{hammond2025risks,anthropic2026multiagent}. These results shift attention from individual model capability toward the dynamics generated by interaction itself. Much of this literature nevertheless studies communication through direct agent-to-agent exchange, structured teams, or explicit deliberative rounds. Recent agent systems point toward a different interaction regime: Moltbook provides an agent-native social environment organized around persistent public activity \citep{jiang2026moltbook}, while the OpenAI--Hugging Face incident showed nominally isolated agents discovering an unsanctioned shared message board and using it to coordinate collective activity \citep{greenblatt2026huggingface}. Such environments have no natural global turn in which every agent observes and responds to every other agent. We therefore study an asynchronous shared communication substrate in which one focal agent becomes active at a time, observes only a bounded portion of the evolving public state, and may change what subsequent agents encounter.

\textbf{Statistical physics of collective dynamics and control.}
Random sequential interaction, in turn, has a long history in statistical-physics models of collective behavior, including opinion dynamics, spreading processes, and convention formation \citep{castellano2009statistical}. Naming-game models show how repeated local interactions can generate population-wide conventions and how network structure shapes their propagation \citep{baronchelli2006sharp,dallasta2006nonequilibrium}, while recent work has begun connecting statistical-mechanical descriptions directly to interacting LLM populations \citep{ashery2025emergent,el2026physics}. A parallel literature studies how collective dynamics can be deliberately steered. Committed agents can act as zealots or pinning controls \citep{masuda2015opinion}, while network-control approaches characterize which components must be actuated, how constrained perturbations propagate, and what resources are required to drive a system toward a desired state \citep{yan2012energy,cornelius2013realistic,ruths2014control,loiudice2015structural,liu2016control}; more recent work considers structured intervention in information diffusion and collective LLM beliefs \citep{chen2026community,he2026belief}. Yet these approaches generally represent communication through local contacts, network links, or propagation events rather than through an evolving semantic message board that agents sample asynchronously. Consequently, the feedback-control problem considered here remains comparatively unexplored: a controller observes only a finite sample of the population, intervenes through the same shared communication substrate under a bounded actuation budget, and is evaluated not only by whether it moves the population but by how efficiently limited sensing and actuation are converted into directed collective response.
